\documentclass[12pt]{article}
\usepackage[utf8]{inputenc}
\usepackage[T1]{fontenc}
\usepackage[margin=1in]{geometry}
\usepackage{setspace}
\usepackage{xcolor}
\usepackage[colorlinks=true, urlcolor=blue!60!black, linkcolor=blue!60!black]{hyperref}
\setlength{\parindent}{0em}
\setlength{\parskip}{0.45em}
\renewcommand{\baselinestretch}{1.0}
\begin{document}
\begin{center}
{\Large \textit{Zelman v. Simmons-Harris}, 536 U.S. 639 (2002)\par}
\vspace{0.3em}
{\normalsize Case Brief --- One-Page Edition\par}
\end{center}

\textbf{Memory Jogger.} Cleveland's vouchers survived: neutral aid that reaches religious schools only through parents' genuine and independent choice is constitutional (\href{https://supreme.justia.com/cases/federal/us/536/639/}{Justia}).

\textbf{Facts.} With Cleveland's failing schools under state control, Ohio gave low-income families tuition scholarships for private schools, religious or secular, or adjacent public schools, and tutoring grants for public-school students. Of 56 participating private schools, 46 were religious, and 96\% of scholarship students attended them; community and magnet schools were also available.

\textbf{Procedural History.} Taxpayers prevailed in the district court and a divided Sixth Circuit.

\textbf{Votes.} Reversed, 5--4: Rehnquist for the Court; O'Connor and Thomas concurred; Stevens, Souter, and Breyer dissented.

\textbf{Holding.} A neutral program aiding a broad class, whose members direct aid to religious schools wholly by their own private choice, does not violate the Establishment Clause.

\textbf{Reasoning.} Under \textit{Mueller}, \textit{Witters}, and \textit{Zobrest}, true private choice breaks the circuit between government and religion. The program was neutral, even financially disfavored religious schools, and offered genuine secular options once community and magnet schools are counted. Souter charged ``verbal formalism'' and a Hobson's choice; Stevens and Breyer warned of religious strife. \emph{Later:} \textit{Espinoza} and \textit{Carson v. Makin} require including religious schools once a State funds private schooling.

\textbf{Examples.} \emph{Same-Side:} savings accounts flowing entirely to religious academies. \emph{Opposite-Side:} vouchers limited to one denomination. \emph{Fence-Sitter:} faith-based rehabilitation programs with a waitlisted secular option.

\textbf{Critique.} The outcome turned on the denominator of choices, and States now face an all-or-nothing choice on private-school aid.

\textbf{Key Quotation.} ``[W]here a government aid program is neutral with respect to religion, and provides assistance directly to a broad class of citizens who, in turn, direct government aid to religious schools wholly as a result of their own genuine and independent private choice, the program is not readily subject to challenge under the Establishment Clause.'' --- Chief Justice Rehnquist, 536 U.S. at 652.

\end{document}
